\section{Pruebas relacionadas con varianzas}

La docimasia de varianzas y dispersiones es indispensable no solo para evaluar uniformidad o riesgo, sino para validar los supuestos de homoscedasticidad en modelos paramétricos.

\begin{teorema}[Prueba $\chi^2$ para la varianza de una población normal]
	Para contrastar si la varianza de una población normal es igual a un valor hipotético $H_0: \sigma^2 = \sigma_0^2$, calculamos la varianza muestral $S^2$ de una muestra de tamaño $n$. El estadístico de prueba es:
	\begin{align}
		\chi^2 = \frac{(n-1)S^2}{\sigma_0^2}.
	\end{align}
	Bajo $H_0$, este estadístico sigue una distribución chi-cuadrada con $\nu = n - 1$ grados de libertad.
	Para una prueba de dos colas ($H_a: \sigma^2 \neq \sigma_0^2$), rechazamos $H_0$ al nivel $\alpha$ si $\chi^2 < \chi^2_{1-\alpha/2, n-1}$ o si $\chi^2 > \chi^2_{\alpha/2, n-1}$.
\end{teorema}

\begin{teorema}[Prueba $F$ para la comparación de dos varianzas poblacionales]
	Para contrastar la hipótesis de igualdad de varianzas entre dos poblaciones normales independientes ($H_0: \sigma_1^2 = \sigma_2^2$), extraemos muestras de tamaños $n_1$ y $n_2$ y calculamos sus varianzas $S_1^2$ y $S_2^2$. El estadístico de prueba es el cociente de las varianzas muestrales:
	\begin{align}
		F = \frac{S_1^2}{S_2^2}.
	\end{align}
	Bajo $H_0$, este estadístico sigue una distribución $F$ de Snedecor con $\nu_1 = n_1 - 1$ grados de libertad en el numerador y $\nu_2 = n_2 - 1$ en el denominador.
	Para una prueba de dos colas ($H_a: \sigma_1^2 \neq \sigma_2^2$), se rechaza $H_0$ si $F > F_{\alpha/2, \nu_1, \nu_2}$ o si $F < F_{1-\alpha/2, \nu_1, \nu_2} = \frac{1}{F_{\alpha/2, \nu_2, \nu_1}}$.
\end{teorema}
